\section{Pendahuluan}

\subsection{Latar Belakang: Kegagalan Model Klasik dalam Krisis Nyata}

Dalam ekonomi neoklasik, sering diasumsikan bahwa pelaku pasar bertindak rasional (\textit{Rational Choice Theory}) dan pasar efisien (\textit{Efficient Market Hypothesis}). Namun, realitas sosial-ekonomi sering menunjukkan sebaliknya:
\begin{enumerate}
    \item \textbf{Irasionalitas Pasar:} Fenomena \textit{panic selling}, \textit{herding behavior} (ikut-ikutan), dan \textit{FOMO} (Fear of Missing Out) yang sering memicu \textit{bubble} atau \textit{crash} finansial tidak dapat dijelaskan sepenuhnya oleh probabilitas klasik (Kolmogorovian).
    \item \textbf{Ketidakpastian Ekstrem:} Model Black-Scholes standar mengasumsikan volatilitas konstan dan distribusi normal, padahal data riil menunjukkan ``ekor gemuk'' (\textit{fat tails}) dan memori jangka panjang (\textit{long-memory effects}), seperti yang dijelaskan dalam teori Fraktal.
    \item \textbf{Dilema Kebijakan:} Pembuat kebijakan sering menghadapi masalah inkonsistensi waktu (\textit{time inconsistency}), di mana kebijakan yang optimal hari ini menjadi tidak kredibel di masa depan karena perubahan ekspektasi publik.
\end{enumerate}

\subsection{Studi Kasus Kajian: Stabilitas Keuangan dan Kebijakan Moneter di Pasar Negara Berkembang (Emerging Markets)}

Studi kasus secara spesifik dapat diterapkan pada \textbf{Pasar Saham Indonesia (IHSG)} dan kebijakan \textbf{Bank Indonesia}, mengingat karakteristik pasar berkembang yang memiliki volatilitas tinggi dan kerentanan terhadap sentimen global.

\subsubsection{Penerapan Teori pada Kasus}
Berikut adalah bagaimana setiap elemen teori diterapkan untuk membedah kasus ini:

\begin{enumerate}
    \item \textbf{Memodelkan Psikologi Investor yang ``Tidak Rasional'' (Quantum Cognition)}
    \begin{itemize}
        \item \textit{Masalah:} Mengapa investor ritel di Indonesia sering melakukan \textit{panic selling} saat ada berita buruk sedikit saja, melanggar prinsip \textit{Sure Thing Principle}?
        \item \textit{Solusi Teori:} Menggunakan \textbf{Quantum-like Influence Diagrams} dan \textbf{Quantum Interference}. Keputusan investor tidak diperlakukan sebagai biner (jual/beli) yang deterministik, tetapi sebagai superposisi state. Interferensi kuantum (konstruktif/destruktif) digunakan untuk memodelkan bagaimana bias kognitif dan ketakutan ``mengganggu'' probabilitas keputusan rasional.
    \end{itemize}

    \item \textbf{Mengukur Risiko Pasar yang ``Kasar'' (Fractal Markets)}
    \begin{itemize}
        \item \textit{Masalah:} Prediksi risiko menggunakan model Gaussian sering meremehkan kemungkinan krisis (Black Swan events).
        \item \textit{Solusi Teori:} Menerapkan \textbf{Fractal Dimensions} dan menghitung \textbf{Hurst Exponent} pada data runtun waktu IHSG. Ini akan membuktikan bahwa pasar tidak efisien (menolak EMH) dan mengikuti \textbf{Fractal Market Hypothesis}, sehingga manajemen risiko harus disesuaikan menggunakan \textbf{Fractional Black-Scholes Models}.
    \end{itemize}

    \item \textbf{Strategi Kebijakan Moneter yang Kredibel (Quantum Game Theory)}
    \begin{itemize}
        \item \textit{Masalah:} Bank Indonesia sering menghadapi dilema antara menaikkan suku bunga untuk menahan inflasi atau menurunkannya untuk memacu pertumbuhan (Dilema Barro-Gordon). Publik sering tidak percaya komitmen bank sentral.
        \item \textit{Solusi Teori:} Memodelkan interaksi ini sebagai \textbf{Quantum Barro-Gordon Game}. Dengan menggunakan strategi kuantum (entanglement strategi kebijakan), dicari \textbf{Nash Equilibrium} baru yang dapat memecahkan masalah inkonsistensi waktu dan mengurangi bias inflasi, menciptakan kebijakan yang lebih dipercaya publik.
    \end{itemize}

    \item \textbf{Prediksi dan Optimasi Portfolio Cepat (Computational Methods)}
    \begin{itemize}
        \item \textit{Masalah:} Kebutuhan prediksi \textit{real-time} untuk \textit{High-Frequency Trading} atau pengawasan pasar.
        \item \textit{Solusi Teori:} Menggunakan \textbf{Quantum-inspired Neural Networks (Qubit/Qutrit)} untuk memprediksi pergerakan harga saham dengan akurasi lebih tinggi daripada Neural Network klasik. Selain itu, menggunakan \textbf{Statistical Arbitrage} berbasis kointegrasi untuk menemukan peluang keuntungan bebas risiko di antara saham-saham LQ45.
    \end{itemize}
\end{enumerate}

\subsection{Tujuan Penelitian}

\begin{enumerate}
    \item Membangun model matematika berbasis formalisme kuantum untuk menjelaskan anomali perilaku investor di pasar saham.
    \item Mengevaluasi tingkat efisiensi dan kekasaran (\textit{roughness}) pasar saham menggunakan dimensi fraktal.
    \item Merumuskan strategi kebijakan moneter optimal menggunakan pendekatan Teori Permainan Kuantum untuk memitigasi inflasi.
    \item Mengimplementasikan algoritma komputasi hibrida (Neural Networks \& Persamaan Diferensial Stokastik) untuk simulasi prediksi pasar.
\end{enumerate}

\subsection{Urgensi Penelitian}
Penelitian ini menunjukkan bahwa alat-alat matematis canggih di Fisika (Ruang Hilbert, Operator Hermisian, Integral Lintasan) bukan hanya abstrak, tetapi merupakan alat paling ampuh untuk memodelkan kompleksitas sistem sosial-ekonomi yang tidak dapat diselesaikan oleh ilmu ekonomi konvensional.